% Gerado por scripts/06_dissertacao.py a partir de output/modelos e output/tabelas.
% Não editar à mão: a próxima execução sobrescreve este arquivo.
\begin{table}[htb]
  \caption{Escala dos domínios de estimação}
  \label{tab:escala}
  \centering
  \espacosimples\idpcorpodez
  \begin{tabular}{@{}lrrrr@{}}
    \toprule
    \textbf{Domínio (desfecho)}                & \textbf{Pessoas-} & \textbf{Células} & \textbf{Ocupa-} & \textbf{UPAs} \\
                                               & \textbf{transições} & \textbf{estimadas} & \textbf{ções} & \\
    \midrule
    Todas as origens (pareamento)              & 3.963.359    & 3.960.584    & 413        & 31.145 \\
    Pares (saída do emprego)                   & 3.178.582    & 3.176.283    & 413        & 31.047 \\
    Destino ocupado (muda de ocupação)         & 2.882.225    & 2.880.049    & 413        & 31.036 \\
    AIOE nos dois lados (margens direcionais)  & 2.879.139    & 2.876.964    & 413        & 31.036 \\
    Origem formal (informalização)             & 1.677.310    & 1.676.862    & 411        & 30.360 \\
    \bottomrule
  \end{tabular}
  \fonte{Elaboração própria a partir da PNADC/IBGE.}
  \nota{Origens de 2019T1 a 2025T3. Pessoas-transições não são pessoas únicas: o
        mesmo indivíduo contribui com até quatro pares. A coluna de células é o
        número de observações efetivamente estimadas em cada modelo.}
\end{table}
